\begin{ans}{1.4.1}
  (1) 準同型 $\phi \colon \mathbb{Z}[x] \to \mathbb{Z}[\sqrt{2}]$ を $\phi(x) = \sqrt{2}$ を
  満たすように定めると, $\phi$ は同型
  $\mathbb{Z}[x]/(x^2 - 2) \cong \mathbb{Z}[\sqrt{2}]$ をひきおこす.
  $\mathbb{Z}[\sqrt{2}]$ のイデアル $(3 + \sqrt{2})$ は, この同型により,
  $\mathbb{Z}[x]/(x^2 - 2)$ において $x + 3 + (x^2 - 2)$ が生成するイデアル
  $(x + 3,\ x^2 - 2)/(x^2 - 2)$ に対応するので,
  \begin{align*}
    \mathbb{Z}[\sqrt{2}]/(3 + \sqrt{2})
    &\cong \nestquot{\mathbb{Z}[x]/(x^2 - 2)}{(x + 3,\ x^2 - 2)/(x^2 - 2)} \\
    &\cong \mathbb{Z}[x]/(x + 3,\ x^2 - 2) \\
    &\cong \nestquot{\mathbb{Z}[x]/(x + 3)}{(x + 3,\ x^2 - 2)/(x + 3)}.
  \end{align*}
  ここで準同型 $\psi \colon \mathbb{Z}[x] \to \mathbb{Z}$ を $\psi(x) = -3$ を満たすように
  定めると, $\psi$ は同型 $\mathbb{Z}[x]/(x + 3) \cong \mathbb{Z}$ をひきおこす.
  $\mathbb{Z}[x]/(x + 3)$ において $x^2 - 2 + (x + 3)$ が生成するイデアル
  $(x + 3,\ x^2 - 2)/(x + 3)$ は, この同型により, $\mathbb{Z}$ において
  $\psi(x^2 - 2) = 7$ が生成するイデアル $7\mathbb{Z}$ に対応するので, 結局
  \[
    \mathbb{Z}[\sqrt{2}]/(3 + \sqrt{2}) \cong \mathbb{Z}/7\mathbb{Z} = \mathbb{F}_7
  \]
  である.

  \medskip
  (2) (1) と同様にして
  \[
    \mathbb{Z}[\sqrt{5}]/(1 + 2\sqrt{5}) \cong \mathbb{Z}[x]/(2x + 1,\ x^2 - 5).
  \]
  ここで
  \begin{align*}
    19 &= -4(x^2 - 5) + (2x - 1)(2x + 1), \\
    x + 10 &= 10(2x + 1) - x \cdot 19
  \end{align*}
  より,
  \[
    (x + 10,\ 19) \subset (2x + 1,\ x^2 - 5).
  \]
  また
  \begin{align*}
    2x + 1 &= 2(x + 10) - 1 \cdot 19, \\
    x^2 - 5 &= (x - 10)(x + 10) + 5 \cdot 19
  \end{align*}
  より逆の包含関係も成り立ち,
  \[
    (x + 10,\ 19) = (2x + 1,\ x^2 - 5)
  \]
  である. よって
  \[
    \mathbb{Z}[\sqrt{5}]/(1 + 2\sqrt{5}) \cong \mathbb{Z}[x]/(x + 10,\ 19).
  \]
  あとは
  \[
    \nestquot{\mathbb{Z}[x]/(x + 10)}{(x + 10,\ 19)/(x + 10)}
  \]
  に対して (1) と同様にするか,
  \[
    \nestquot{\mathbb{Z}[x]/(19)}{(x + 10,\ 19)/(19)}
  \]
  に対して命題 1.4.14 を適用 ($A = \mathbb{Z}$, $I = 19\mathbb{Z}$)
  したのちに例 1.4.11 を適用 ($A = \mathbb{Z}/19\mathbb{Z}$) すれば
  $\cong \mathbb{F}_{19}$ がわかる.
\end{ans}

\begin{ans}{1.4.2}
  $A = \mathbb{C}[x]/(x^2)$, $B = \mathbb{C}$.
\end{ans}

